\documentclass[10pt,a4paper]{article}
\usepackage[margin=23mm]{geometry}
\usepackage[T1]{fontenc}
\usepackage{lmodern,microtype,amsmath,amssymb,amsthm,mathtools}
\usepackage{booktabs,longtable,array,enumitem,xcolor,graphicx}
\usepackage[colorlinks=true,linkcolor=blue!55!black,urlcolor=blue!55!black,citecolor=blue!55!black]{hyperref}
\usepackage{fancyhdr}
\pagestyle{fancy}\fancyhf{}\fancyhead[L]{Acquiring the first-event law}
\fancyhead[R]{5 November 2026}\fancyfoot[C]{\thepage}
\setlength{\headheight}{14pt}
\setlength{\parskip}{4pt}\setlength{\parindent}{0pt}
\setlist{itemsep=2pt,topsep=4pt}
\newtheorem{theorem}{Theorem}
\newtheorem{lemma}{Lemma}
\newtheorem{corollary}{Corollary}
\newcommand{\TV}{d_{\rm TV}}
\newcommand{\E}{\mathbb E}
\newcommand{\Prb}{\mathbb P}
\newcommand{\one}{\mathbf 1}
\newcommand{\Law}{\mathcal L}
\newcommand{\ceil}[1]{\left\lceil#1\right\rceil}
\newcommand{\eps}{\varepsilon}
\newcommand{\sig}{\sigma}
\DeclareMathOperator{\supp}{supp}
\hypersetup{pdftitle={Acquiring the first-event law},pdfsubject={Local experimental access, joint marked-path uncertainty, finite generation and intervention decisions}}
\title{\textbf{Acquiring the first-event law}\\[4pt]
\large Local experiments, marked-path uncertainty, and decisions before rare failure}
\author{Research continuation with preserved governing sources}
\date{5 November 2026}
\begin{document}
\maketitle
\begin{abstract}
We construct a finite interface from local experimental evidence to the
joint first-disconnection time, terminal state and a requested recent
precursor window in a declared interacting Gibbs reliability model.
The abstract uncertainty bridge applies more generally whenever its
finite-chain interfaces are supplied. The first route acquires a reusable
model and certifies its full
event-record law. Its central new bound compares short marked-Palm
proposals in relative entropy, controls their accepted submeasures in
Hellinger distance, and separately controls the geometric waiting clock.
The resulting uncertainty budget has no inverse global-rarity or mean
lifetime factor. The second route uses controlled local experiments as
probability coins. Balanced-field sampling supplies true static updates;
censored release races and a finite Bernoulli factory supply true operating
updates without fitting a fine kinetic table. An explicit weight-ratio
construction closes the unknown-static-weight interface.

Both routes compose with the completed finite first-event sampler and a
paid safe preparation, under a certified safe-mass floor, to give
bounded-loss guarantees for a finite menu
of constant interventions. They retain their different access and costs:
the first pays fine acquisition once, while the second consumes fresh
physical queries for later forecasts. Model closure, faithful copies,
calibrated fields, priced pins, timing, mixing and static contrast margins
remain substantive. The result is a mathematical and experimental-interface
construction, not an economical prototype. New executions cover finite
primitives and certificate arithmetic, not a general large Gibbs engine
or a physical apparatus. Independent foundational and mechanical
constructions, all earlier originals and useful failed attempts are
preserved in the companion package.
\end{abstract}

\section{The question and the completed consequence}
Can local controlled experiments support a forecast of \emph{what fails,
when it fails, and the recent sequence leading to it}, without waiting
for a global failure? Can that forecast choose an actual permitted
intervention with a quantitative loss guarantee?

The completed consequence is conditional but stronger than a CDF or mean
certificate. For all-terminal network disconnection under the declared
interacting local Gibbs and clock model, the
experimental interfaces described below construct a finite generator for
the entire requested event record. With specified confidence over
acquisition and shared calibration, its output law is within a specified
total-variation error of the physical mission law. Every bounded loss
of that record is consequently controlled. A finite menu can be compared
before deploying the chosen intervention and performing its finite safe
preparation. There is no unsuccessful full-lifetime simulation factor.

The two realizations are deliberately distinct. A fitted-model generator
can be used offline after its fine parameter acquisition; its robustness
is proved in Section~\ref{sec:bridge}. An experiment-fed generator uses
fresh physical randomness and does not export a finely estimated rate
table; its detailed construction is in the separate complete proof and
the operational derivation below. Neither result asserts that finite
local data discover omitted external mechanisms.

\subsection{Why the earlier CDF theorem was insufficient}
Consider states $a,b_1,b_2$, start at $a$, and let the event be
$B=\{b_1,b_2\}$. Set
\[
Q(a,b_1)=\lambda r,\quad Q(a,b_2)=\lambda(1-r),\quad
Q(b_1,a)=Q(b_2,a)=1.
\]
This is a full-support reversible chain. The first-event CDF is
$1-e^{-\lambda t}$ for every $r\in(0,1)$, while its terminal law is
$(r,1-r)$. Replacing $r=0.99$ by $r=0.01$ preserves the entire CDF
exactly and changes terminal-state TV by $0.98$. Thus even perfect CDF
agreement does not supply the missing marked-law interface. The earlier
relative whole-CDF operation remains valuable for very small deadline
probabilities and is not replaced by the present absolute-TV guarantee.

\subsection{The retained record and the paid information}
Let $Q$ be a finite irreducible CTMC, $B$ a nonempty proper event set,
$\pi$ its stationary law, and $p=\pi(B)$. A succinct state need not
enumerate the entire state space. Choose a common lazy uniformizer
$R\ge2\max_x[-Q(x,x)]$, a certified normalized additive-symmetrization
gap $g>0$ and a stationary mass floor $\zeta>0$. Write
$P=I+Q/R$. The retained record is
\[
\mathcal R_{H,s}=
\bigl(\lfloor2^sR\tau\rfloor,\ X_\tau,
\text{state and ordered changes in the last }H\text{ normalized units}\bigr),
\]
with precursor timestamps quantized at the same mesh and history
truncated at the mission start. Null uniformization updates are not
printed. A physical window $h$ pays $H=\lceil Rh\rceil$.
The complete uneventful past is not retained. Its output length would
itself carry a lifetime cost.

All TV comparisons concern this same finite record. A finite-valued
timestamp is not claimed TV-close to an unrounded continuous time.
The accepted finite-bit compiler pays clock digits, update caps,
conditional-state error, rejection caps, retained history and memory.

\subsection{Precisely what is reused}
The first-event sampler, finite clock, acquisition/preparation and
correlated static constructions at fixed source revision
\texttt{719e201332992002528c221c49f80f69431485f6} are accepted machinery
\cite{governing}. In the rare branch their central identity is
\begin{equation}
\frac{d\mathsf P_L}{d\Law_\pi(X_{[0,L]})}
=\frac{N_B}{p(1+L)},\qquad
w=\frac{\one_{\{H<\tau\le L\}}}{N_B}.
\label{eq:palmBaseline}
\end{equation}
The proposal starts from a rare stationary state placed either at time
zero or at a uniformly marked potential event, then evolves forward and
backward using $P,P^*$. Its acceptance mean is
$\mu=q_G/[p(1+L)]$. For $L\ge16(H+1)/g$ and $Lp\le1$,
$\mu\ge g/8$. This is the nonrare acceptance floor that the new bridge
uses. Long waiting is represented by a calibrated geometric integer,
not by executing all intervening updates. Buffers handle dependence
and arbitrary supplied initial states at their proved structural cost.

The correlated static solver still requires its declared factor graph,
certified pinned sampling/mixing and fixed conditional log-surprisal
contrast. The new experiment-fed construction implements these static
queries from local experimental access; it does not remove the margins.
The original code and conclusions are preserved, rather than edited to
fit the new theorem.

\section{A quantitative bridge from model uncertainty to event records}
\label{sec:bridge}

This section constructs the missing marked-path comparison. It does not
deduce it from the older CDF comparison. The two comparisons have different
targets and remain distinct.

\subsection{Two-model input and the retained record}
Let $Q,\widehat Q$ be finite irreducible generators on the same encoded
state space, with invariant laws $\pi,\widehat\pi$ and the same nonempty
proper event set $B$. Use a common rational
$R\ge2\max\{\max_x[-Q(x,x)],\max_x[-\widehat Q(x,x)]\}$ and put
$P=I+Q/R$, $\widehat P=I+\widehat Q/R$.
Assume both normalized additive symmetrizations have gap at least $g>0$
and both invariant laws have mass floor $\zeta>0$.
Stationary reverses are denoted by stars. Common transition support is a
premise for the finite divergence envelopes below.

The event record has the same meaning as in the accepted first-event
construction: the quantized first time, terminal state and ordered changes
in the last $H$ normalized time units, truncated at the supplied initial
time. The comparison first proves a statement for continuous records;
applying any common finite timestamp quantization cannot increase TV.

Define finite bounds
\begin{align}
\chi_B&\ge D_{\rm KL}(\pi(\cdot\mid B)\Vert\widehat\pi(\cdot\mid B)),\\
\kappa_+&\ge\sup_xD_{\rm KL}(P(x,\cdot)\Vert\widehat P(x,\cdot)),\\
\kappa_-&\ge\sup_xD_{\rm KL}(P^*(x,\cdot)\Vert\widehat P^*(x,\cdot)),\\
d_p&\ge|\log(\pi(B)/\widehat\pi(B))|.
\end{align}
These are model-confidence bounds, not assumed observations of a rare
trajectory. A local experimental construction of them is given below.

\subsection{Square-root distance of marked proposals}
Write
$\mathsf h^2(\nu,\omega)=\int(\sqrt{d\nu}-\sqrt{d\omega})^2$
also for subprobability measures. For probabilities,
$\TV(\nu,\omega)\le\mathsf h(\nu,\omega)$ and
$\mathsf h^2(\nu,\omega)\le D_{\rm KL}(\nu\Vert\omega)$.

Use the accepted length-$L$ marked-Palm proposal: its initial-mark branch
has probability $1/(1+L)$ and its internal-mark branch has probability
$L/(1+L)$. The two models use the same branch, Poisson count, marked
position and uniform timestamps; only their conditional rare state and
forward/reverse updates differ.

\begin{lemma}[Proposal information bound]\label{lem:palmKL}
For the two marked proposals $\mathsf P_L,\widehat{\mathsf P}_L$,
\begin{equation}
D_{\rm KL}(\mathsf P_L\Vert\widehat{\mathsf P}_L)
\le D_L:=
\chi_B+\frac{L(L+2)}{2(L+1)}(\kappa_++\kappa_-).
\label{eq:DL}
\end{equation}
\end{lemma}
\begin{proof}
Retain the branch and marked position as auxiliary variables. The relative
entropy chain rule first pays $\chi_B$ at the marked state and then the
appropriate row divergences for all generated updates. The auxiliary
count/time laws are identical. In the internal-mark branch, conditional
on $n=M+1$ events, the marked position is uniform in $1,\ldots,n$.
Its expected backward count is $(n+1)/2$ and its expected forward count
is $(n-1)/2$. Since $\mathbb E M=L$, weighting the two branches gives
expected forward and backward counts each equal to
$L(L+2)/(2(L+1))$. Dropping auxiliary variables can only decrease KL.
\end{proof}

Let $G$ mean that the first hit within the block lies in $(H,L]$, and
$N_B$ count its potential-event states in $B$, including time zero and
null updates. Both proposals have the same measurable acceptance
function $w=\mathbf 1_G/N_B$. Let their acceptance means be
$\mu,\widehat\mu$.

\begin{lemma}[Accepted-submeasure comparison]\label{lem:acceptKL}
If $\mu,\widehat\mu\ge q_0>0$ and $z=\sqrt{D_L/q_0}$, then
\begin{equation}
\TV(\mathsf P_L(\cdot\mid\mathrm{accept}),
     \widehat{\mathsf P}_L(\cdot\mid\mathrm{accept}))\le z,
\qquad
|\log(\mu/\widehat\mu)|\le2z.
\label{eq:acceptKL}
\end{equation}
\end{lemma}
\begin{proof}
Multiplication by $0\le w\le1$ cannot increase square-root distance.
For the accepted submeasures $a=w\mathsf P_L$ and
$\widehat a=w\widehat{\mathsf P}_L$,
\begin{equation}
\mathsf h^2(a,\widehat a)=
(\sqrt\mu-\sqrt{\widehat\mu})^2+
\sqrt{\mu\widehat\mu}\,
\mathsf h^2(a/\mu,\widehat a/\widehat\mu).
\label{eq:submass}
\end{equation}
The left side is at most $D_L$. The second term gives the TV bound.
For the mass bound, use
$|\log\sqrt\mu-\log\sqrt{\widehat\mu}|
\le|\sqrt\mu-\sqrt{\widehat\mu}|/\sqrt{q_0}$.
\end{proof}

The acceptance denominator is $q_0$, not $\pi(B)$. Under the accepted
uniform block conditions it can be taken as $g/8$.
Because the stationary block success probability is
$q_G=p(1+L)\mu$, the two geometric parameters obey
\begin{equation}
|\log(q_G/\widehat q_G)|\le d_p+2z.
\end{equation}
The Bernoulli coupling bound therefore gives
\begin{equation}
\TV(\operatorname{Geom}_0(q_G),
    \operatorname{Geom}_0(\widehat q_G))
\le\min\{1,e^{d_p+2z}-1\}.
\label{eq:geom-model}
\end{equation}
The product of this waiting index and an independent accepted path
differs by at most the sum of \eqref{eq:acceptKL} and
\eqref{eq:geom-model}.

\subsection{Complete event-law comparison, including the initial state}
For an integer buffer $b$, let $\eta=e^{-gb}/\zeta<1$.
For each model separately let
\begin{equation}
S_Q=\frac{\eta}{1-\eta}
+\frac{1+\eta}{1-\eta}\frac{q_Q(H)+q_Q(b)}{q_{G,Q}}
+(1+\eta)q_Q(b),
\label{eq:SQ}
\end{equation}
and define $S_{\widehat Q}$ analogously. These are the full structural
bounds of the accepted block construction, including its joint
pre-block buffer charge and unnormalized initial survivor law.

\begin{theorem}[Joint first-event model stability]\label{thm:stability}
Suppose both models satisfy the block conditions and have acceptance
means at least $q_0$. For every common initial law $\nu$,
\begin{equation}
\TV(\Law_{\nu,Q}(\mathcal R_{H,s}),
    \Law_{\nu,\widehat Q}(\mathcal R_{H,s}))
\le
S_Q+S_{\widehat Q}
+\sqrt{b\kappa_+/2}
+z+\min\{1,e^{d_p+2z}-1\}.
\label{eq:master}
\end{equation}
An initial-law TV discrepancy adds its value. There is no inverse
stationary-rarity or mean-lifetime factor.
\end{theorem}
\begin{proof}
Compare each true event record with its own ideal block construction.
Those comparisons cost $S_Q$ and $S_{\widehat Q}$. Couple the observed
initial length-$b$ paths. Under Poisson uniformization their relative
entropy is at most $b\kappa_+$, since they have the same initial law;
Pinsker gives the displayed prefix cost. On a matching early hit both
algorithms return the same early record. On matching survival their late
outputs consist of the independent geometric index and successful path,
shifted by the same integer calendar offset. The preceding two lemmas
control that product law. All bounds are uniform in $\nu$. Coupling
initial states adds an initial-law error by contraction of a Markov
channel, and final quantization cannot increase TV.
\end{proof}

\subsection{Explicit envelopes acquired from local factors and rates}
Suppose $W,\widehat W$ are unnormalized Gibbs weights satisfying
\begin{equation}
|\log\widehat W(x)-\log W(x)|\le E
\quad\text{for every }x,
\qquad
|\log(\widehat Q(x,y)/Q(x,y))|\le a
\quad(x\ne y).
\label{eq:local-envelope}
\end{equation}
For the reversible true and learned local models in the accepted active
learner, both reverse kernels equal the corresponding forward kernel.

\begin{lemma}[Finite local-to-proposal bound]\label{lem:localKL}
Under these reversible premises and the common lazy uniformizer,
\begin{equation}
d_p=2E,\qquad
\chi_B=2E^2,\qquad
\kappa_+=\kappa_-=(e^a-1)^2
\label{eq:envelopes}
\end{equation}
are valid conservative bounds. In particular,
\begin{equation}
D_L\le 2E^2+(L+1)(e^a-1)^2.
\label{eq:reversibleDL}
\end{equation}
\end{lemma}
\begin{proof}
Both the full partition sum and its restriction to $B$ lie within
$e^{\pm E}$ of their learned counterparts, giving $d_p\le2E$.
The conditional density ratio on $B$ lies in $[e^{-2E},e^{2E}]$.
For any two probabilities with ratio in $[e^{-c},e^c]$,
convexity of minus the logarithm gives
$D_{\rm KL}\le c\tanh(c/2)\le c^2/2$; apply $c=2E$.

For a row, put $v=\sum_{y\ne x}\widehat P(x,y)\le1/2$.
The off-diagonal chi-square contribution relative to the learned row is
at most $(e^a-1)^2v$. The diagonal contribution is at most
$(e^a-1)^2v^2/(1-v)\le(e^a-1)^2v$.
Their sum is at most $(e^a-1)^2$. Relative entropy is bounded by
chi-square divergence. Reversibility supplies the same reverse bound,
and $L(L+2)/(L+1)<L+1$.
\end{proof}

For a nonreversible confidence family the same theorem applies with
separately certified reverse-row bounds. Stationary weight ratios and
incoming rates must construct them. The reversible active learner does
not silently acquire arbitrary driven kinetics.


\section{A finite acquired certificate and its intervention meaning}
\label{sec:acquired}

\subsection{A common branch, with no hidden rare-event threshold}
Fix $0<\varepsilon\le1/4$. Conservatively replace a normalized gap
certificate by $\min(g,1)$, and define
\begin{align}
C&=\left\lceil\log_2\frac{128}{\varepsilon\zeta}\right\rceil,
&b&=\lceil C/g\rceil,\\
L&=\left\lceil\frac{1024(H+b+2)}{g\varepsilon}\right\rceil,
&T_{\max}&=\lceil16LC/g\rceil+1.
\label{eq:commonSchedule}
\end{align}
The two models must use these same buffers, block lengths, clock units
and event predicate. Count the learned stationary event probability to a
positive interval $[\ell,u]$, and widen it to
$[p_-,p_+]=[e^{-2E}\ell,e^{2E}u]$ with outward rounding. Refine Count
and require $p_+/p_-\le2$. Under the energy precision below this is a
finite relative-counting task. Both true and learned probabilities then
belong to this interval. All counting failures have their own confidence
budget; this branch rule is used only on the simultaneous good event.

If $p_+\le1/(8L)$, use Theorem~\ref{thm:stability} with $q_0=g/8$.
The accepted structural estimates give $S_Q,S_{\widehat Q}<\varepsilon/32$.
If $p_+>1/(8L)$, then $p_->1/(16L)$. Direct simulation to
$T=\lceil C/(gp_-)\rceil\le T_{\max}$ suffices. For either model,
the killed-semigroup bound
$\Pr_\nu(\tau>T)\le\zeta^{-1/2}e^{-gp_-T}$ is at most
$\varepsilon/128$. Coupling the two finite uniformized trajectories gives
\begin{equation}
\TV(\Law_{\nu,Q}(\mathcal R_{H,s}),
        \Law_{\nu,\widehat Q}(\mathcal R_{H,s}))
\le\sqrt{T\kappa_+/2}+\varepsilon/64.
\label{eq:directStability}
\end{equation}
Thus the direct branch does not reintroduce a numerical $1/p$ cost:
its certified horizon is bounded by $T_{\max}$ before execution.

For example, the following sufficient, deliberately loose precisions
imply a model-law discrepancy below $\varepsilon/4$ in either branch:
\begin{equation}
E\le\frac{\varepsilon\sqrt g}{256},\qquad
a\le\min\left\{1,
\frac{\varepsilon}{256}\sqrt{\frac g{L+1}},
\frac{\varepsilon}{64\sqrt{T_{\max}}}\right\}.
\label{eq:sufficientprecision}
\end{equation}
Indeed $e^a-1\le2a$, and consequently
$D_L\le3\varepsilon^2g/32768$,
$z\le\sqrt3\,\varepsilon/64<\varepsilon/32$,
$d_p\le\varepsilon/128$ and the observed-prefix term is below
$\varepsilon/32$. The geometric term is below $\varepsilon/8$.
The direct bound is smaller. Finite upward rounding can use strict
slack in these inequalities. These are an existence schedule, not an
optimized experimental design and not lower bounds.

\subsection{A certificate that can use an actually measured acceptance}
The uniform schedule is not compulsory. For a chosen $b,L$, suppose
finite calibration provides an outward interval
$\widehat\mu\in[\widehat m_-,\widehat m_+]$ including all numerical and
sampling biases. Equation~\eqref{eq:submass} yields
\begin{equation}
\mu\ge m_-:=
\bigl(\sqrt{\widehat m_-}-\sqrt{D_L}\bigr)_+^2.
\label{eq:observedAcceptance}
\end{equation}
If $m_->0$, use $q_0=\min(m_-,\widehat m_-)$ in the comparison.
For either model its structural bound can be evaluated without knowing
$p$ exactly:
\begin{equation}
S\le\frac{\eta}{1-\eta}
+\frac{1+\eta}{1-\eta}
\frac{H+b+2}{(1+L)m_{\rm low}}
+(1+\eta)p_+(b+1),\qquad \eta=e^{-gb}/\zeta,
\label{eq:instanceStructural}
\end{equation}
where $m_{\rm low}$ is the corresponding acceptance lower bound.
This uses $q(t)\le p(1+t)$ and $q_G=p(1+L)\mu$.
The calibrated proposal must use the intended $G$ and common window;
one cannot insert an unrelated observed acceptance rate. If the
square-root gap closes, this certificate reports ``unresolved'' rather
than a positive floor. Its finite arithmetic and statistical errors are
explicit inputs, not silently erased.

\subsection{Local evidence supplies the uncertainty objects}
Use the accepted reversible, pairwise binary Gibbs acquisition model.
There are $N$ variables, $M$ declared interaction edges, maximum degree
$d$, a finite menu of $K$ constant calibrated interventions $u$, and
full observed model closure. In intervention $u$,
\begin{align}
W^u(x)&=\exp\!\left\{\sum_i(h_i+u_i)x_i+
                         \sum_{ij}J_{ij}x_ix_j\right\},\\
Q^u(x,x^i)&=r_i^u(x_{\partial i})
\sigma\bigl((1-2x_i)(h_i+u_i+\sum_jJ_{ij}x_j)\bigr).
\end{align}
The unknown context clocks have supplied or acquired valid numerical
bounds $0<r_*\le r_i^u(s)\le r^*$. The additive-field model,
local detailed balance, faithful isolated copies, the absence of hidden
memory and the declared graph are premises. Finite observations do not
certify model completeness.

The accepted balanced-field learner, with parameter tolerance $\eta_0$
and kinetic log-rate tolerance $\rho_0$, gives simultaneous envelopes
\begin{equation}
E_u=(N+2M)\eta_0+N(\xi_u+\tau_u),\qquad
a_u=\rho_0+(1+2d)\eta_0+\xi_u+\tau_u.
\label{eq:acquiredEnvelopes}
\end{equation}
Here $\xi_u$ is deployment-field calibration error and $\tau_u$ is
the rational factor/control rounding allowance. Errors in actual timing
and detector outputs must be included in that learner's confidence
certificate; they are not covered by the displayed field allowance.
The original complete acquisition proof and experiment record are
preserved unchanged. In particular, a biased calibration clock is not
assumed to equal an operating clock.

Interval Dobrushin rows, the pinned conditional-probability contrast
condition of the accepted correlated static algorithm, rate bounds and
a mass floor construct common $g,\zeta,R$ for every model in the
retained confidence box. A coarse box may first establish strict
margins; subsequent acquisition refines inside that box. A procedure
with unresolved margins returns that status. It cannot silently turn a
mixing assumption into a learned fact. The static solver is used at
its original graph, contrast and pinned-mixing scope, including its
numerical inverse margins and degree dependence.

Equations~\eqref{eq:acquiredEnvelopes} and
\eqref{eq:sufficientprecision}, or the tighter explicit certificate
\eqref{eq:master}, now specify an acquisition precision sufficient for
the joint event record. This is the new interface. The previous CDF
stability theorem alone supplied none of these marked-path terms.

\subsection{The complete resource bill for the fitted-model route}
Let $D=N+2M$ and $C_{\rm kin}=K\sum_i2^{\deg(i)}$. Omitting only
universal numerical constants, the inherited finite schedules use
\begin{align}
\text{energy reads}\quad&
O\!\left(D\left[\log(2+B_0)+\eta_0^{-2}\right]
                 \log\frac{D(1+\log(2+B_0))}{\delta_E}\right),\\
\text{energy dwell}\quad&
O\!\left(\frac D{r_{\rm cal}}
\left[\log(2+B_0)+\eta_0^{-2}\log(2/\eta_0)\right]
\log\frac{D(1+\log(2+B_0))}{\delta_E}\right),\\
\text{operating releases}\quad&
O\!\left(C_{\rm kin}\rho_0^{-2}
                      \log(C_{\rm kin}/\delta_Q)\right),\\
\text{capped kinetic dwell}\quad&
O\!\left(\frac{C_{\rm kin}}{r_*\rho_0^2}
\log(C_{\rm kin}/\delta_Q)
\log\frac{C_{\rm kin}(1+\rho_0^{-2})}{\delta_Q}\right).
\label{eq:fitCosts}
\end{align}
The original constants, finite censoring and timing construction govern
these summaries. Priced pins and neighbor clamps, context setup, bit
reads, event detector service, field calibration and finite-slew travel
are added at their actual costs. In particular, a pin is not a free
state reset and a detector's clock digits are not a physical response
time guarantee. The degree term and numerical $r_*^{-1}$ can be large.

Once this acquisition is complete, the learned-model sampler pays the
accepted finite static, reverse/forward and clock compilation costs.
Its conservative rare-branch calibration uses
$O(g^{-1}\varepsilon_{\rm sim}^{-2}\log(1/\delta_{\rm sim}))$
length-$L$ Palm proposals; a new draw uses at most
$O(g^{-1}\log(1/\varepsilon_{\rm sim}))$ proposals and a length-$b$
initial segment. Each proposal uses one approximate rare stationary
state and $O(L+\log(1/\theta))$ finite updates, at the prescribed
per-call tolerance $\theta$. Direct-branch updates are bounded by
$T_{\max}$ and the same finite truncation rules. The waiting index is
generated in binary by the preserved finite geometric-clock compiler;
its digits, not its value, are paid. A physical window of duration $h$
requires $H=\lceil Rh\rceil$. Window output and memory scale accordingly.

All formulas depend on local experimental scales, degree, inverse
mixing/static margins, desired accuracy and $\log(1/\zeta)$. There is
no $1/\pi(B)$ or mean-first-event simulation factor. There may still be
exponential costs in degree, physical contrast, mixing margin or
apparatus precision. The earlier thirteen-year serial example is not
made economical by this theorem; its coarse CDF calibration is not a
joint-record certificate.

\subsection{Preparation and decisions that the access really permits}
The abstract comparison is uniform over a common supplied initial law.
For an actual equilibrium-safe mission, use the accepted finite physical
preparation only for menu interventions with the certified bound
$\pi^u(A)\ge1/2$, verified by $p_+^u\le1/2$ from the same widened
count interval. This safe-mass premise is part of the complete operation,
not a consequence of a mixing gap. If this premise is to be certified
only from a two-sided statistical interval, a strict safe-mass slack
(and its required counting precision) is paid; an unresolved boundary
returns that status. In the rare branch the displayed event-mass bound
already supplies ample safe slack. Under the chosen constant intervention,
let the repairable
chain evolve for its certified mixing interval, inspect the state at a
known completion time, retain a safe result, and repeat up to a finite
cap. This requires the physical process to continue through unsafe
states during preparation. It is not an experiment on an irreversibly
destroyed system. Its conditional successful law $\nu_u$ satisfies
\begin{equation}
e^{-\beta_u}\le
\frac{d\nu_u}{d\pi^u(\cdot\mid A)}\le e^{\beta_u},
\qquad A=B^c,
\end{equation}
and its abort probability is at most $\delta_{\rm prep}$. The inherited
inspection, field ramp, mixing, retry, and detector costs are paid. For
the stated safe stationary mass bound, the finite retry cap is logarithmic
in $1/\delta_{\rm prep}$, without a rare-event rejection cost.

The learned safe initial law has safe mass at least $e^{-2E_u}/2$ and
can therefore be sampled with finite capped rejection and prescribed
work. More generally a supplied safe-mass floor $s_*>0$ is possible,
but both physical preparation and safe-law sampling must pay
$O(s_*^{-1}\log(1/\delta_{\rm prep}))$ attempts (and the learned
density-ratio factor). The present constant-cost preparation claim uses
$s_*=1/2$ throughout. The density-ratio bound gives
\begin{equation}
\TV(\nu_u,\widehat\pi^u(\cdot\mid A))
\le \tanh(\beta_u/2)+\tanh(E_u).
\label{eq:prepTV}
\end{equation}
Add the numerical safe-sampling tolerance. Combining preparation,
Theorem~\ref{thm:stability} (or \eqref{eq:directStability}) and the finite
model sampler yields a completely specified error
\begin{equation}
\varepsilon_{{\rm law},u}=
\varepsilon_{{\rm model},u}+\varepsilon_{{\rm sim},u}
+\varepsilon_{{\rm init},u}
+\tanh(\beta_u/2)+\tanh(E_u).
\label{eq:totalAcquiredTV}
\end{equation}
This is a conditional-on-success mission law. For an unconditional
operation, output an explicit preparation-abort symbol and charge its
probability; assigning it loss at most one adds at most
$\delta_{\rm prep}$ to risk.

For every menu intervention, choose any bounded loss
$\ell_u(\mathcal R)\in[0,1]$. It may jointly depend on event time,
terminal damage and the retained precursor window. On the single
simultaneous acquisition/calibration confidence event, the difference
between physical and generated expected loss is at most
$\varepsilon_{{\rm law},u}$. With
\begin{equation}
n\ge\frac{\log(2K/\delta_{\rm dec})}{2\sigma^2}
\label{eq:decisionSamples}
\end{equation}
fresh generated records per intervention, all empirical risks have
additional error at most $\sigma$, with probability at least
$1-\delta_{\rm dec}$. Selecting their minimum gives
\begin{equation}
\operatorname{Risk}(\widehat u)
-\min_{u\in\mathcal U}\operatorname{Risk}(u)
\le2\left(\max_u\varepsilon_{{\rm law},u}+\sigma
+\delta_{\rm prep}\right).
\label{eq:regret}
\end{equation}
Any known intervention cost can be included after scaling total loss
into the declared bounded range. The event-law TV certificate covers
all bounded record losses, but the finite-sample simultaneous statement
above is for the declared $K$ evaluated losses; selecting arbitrary
losses after inspecting the same records needs its own complexity bound.

The decision is made after acquisition and before deploying the chosen
field and starting its preparation. Only the selected system is
physically prepared. This does not grant instantaneous feedback on an
unheld live network: observing a state, computing a choice and ramping
the apparatus take real time. A mission that must remain active during
acquisition or decision requires a separately priced holding, shutdown
or exposure model. The theorem does not supply such access.

Absolute-TV risk guarantees do not give relative accuracy for every
rare damage subtype, nor control unbounded lifetime costs. The acquired
relative whole-CDF theorem remains a separate operation and can be
used alongside the present joint-record guarantee.



% The complete operational construction and execution/comparison sections
% are inserted after the independent experimental proof is finalized.
\section{A second realization: use the experiments as update coins}
The same acquired-event endpoint has a different implementation. A fine
kinetic table can be avoided if new local experiments remain available
during compilation and every forecast. This changes the experimental
bill; it does not turn a finite data set into an unlimited exact simulator.
The following construction was developed independently of the
marked-law bridge. Its complete standalone source and execution record
are preserved in \texttt{next5\_decision}.

Use the same binary pairwise model, with known neighborhoods and unknown
$h_i,J_{ij}$, and the actual menu-dependent rates
\begin{equation}\label{eq:Q}
Q^u(x,x^{(i)})=r_i^u(s)\sig((1-2x_i)(\ell_i(s)+u_i)),
\quad \ell_i(s)=h_i+\sum_jJ_{ij}s_j,
\quad 0<r_*\le r_i^u(s)\le r^*.
\end{equation}
The same $r_i^u(s)$ applies to both values of the focal bit. That is the
declared reversible mechanism. A faithful local cell realizes every
requested neighborhood, the full operating environment, and the actual
focal clock. A held artificial field $z$ has stationary probability
$\sig(\ell_i(s)+z)$ and relaxation rate at least $r_{\rm cal}$, even
though its clock may differ arbitrarily from the operating clock.
Coarse acquisition locates a balanced field with probability in
$[1/4,3/4]$ and certifies the required strict static/mixing margins.

Context realization, additive-field calibration, available slew and
settling, focal state pins, bit reads, event detection and reference timing
have the actual costs
\[
\begin{gathered}
C_{\rm context},\ C_{\rm field}(\xi),\ C_{\rm slew},\ C_{\rm settle},\\
C_{\rm pin},\ C_{\rm read}(\beta),\ C_{\rm event}(\beta),\
C_{\rm timer}(\Delta).
\end{gathered}
\]
Missed exits, dead times, readout error and drift must fit the allocated
conditional error budgets. The operation is finite only if those
services are supplied at the requested precision. The proof does not
build a universal cell, field meter or detector.

For the visible edge-failure bits, coarse intervals must certify
\begin{equation}\label{eq:contrast}
0<q_i^-\le\pi^u(x_i=1\mid x_{-i})\le q_i^+<1,
\qquad q_i^-\ge(q_i^+)^c,
\end{equation}
with fixed $c$, as required by the accepted correlated cut construction.
A pinned Dobrushin margin is $\alpha>0$. If
$|\ell_i(s)|\le B_0$, $|u_i|\le U$, one convenient common certificate is
\begin{equation}\label{eq:gap}
b_E=\lceil2(B_0+U)\rceil,\quad
\zeta=2^{-N(b_E+1)},\quad
\rho=2r^*,\quad R=N\rho,\quad g=\alpha r_*/R.
\end{equation}
Every complete state has mass at least $\zeta$: the bound on a full
conditional bit probability also bounds partial conditionals by the
tower property, and these factor the state probability. Heat-bath
Dirichlet-form domination by $r_*$ supplies the normalized gap.

There are two explicit coarse acquisition options. Enumerating all
local contexts costs $\sum_i2^{\deg(i)}$ and is polynomial only at
fixed degree. In the pairwise model, the accepted $D=N+2M$ energy-context
learner instead reconstructs all local logits to error
$(1+2d)\eta_0$. Choosing
$\eta_0\le1/[8(1+2d)]$ and field error at most $1/8$ makes
$z=-\operatorname{clip}_{[-B_0,B_0]}(\widehat\ell_i(s))$ balance every
later context without exceeding the declared field range. The sufficient
interval row test
$\sum_j\tanh((|\widehat J_{ij}|+2\eta_0)/4)<1$
certifies a pinned Dobrushin margin when it has strict slack. Interval
sums of positive and negative couplings bound every conditional logit.
This second option avoids scanning an arbitrary clock table, but it
requires that the stated sufficient row certificate succeed. Refinement
and unresolved-premise outcomes are explicit. Neither option removes
the fine field metrology needed in the update construction below.

\section{A physical true-Gibbs update without fine energy estimation}

\begin{lemma}[Balanced two-coin construction]\label{lem:gibbs}
Suppose a balanced physical coin has success probability
$q=\sig(\ell+z)\in[1/4,3/4]$. For $a=e^{u-z}$ the following experiment has
output law $\operatorname{Bernoulli}(\sig(\ell+u))$: choose branch 1 with
probability $a/(1+a)$ and branch 0 otherwise; toss a $q$ coin; on branch 1
accept a success and return 1; on branch 0 accept a failure and return 0;
otherwise repeat. The acceptance probability on each attempt is at least
$1/4$.
\end{lemma}
\begin{proof}
The probabilities of accepted 1 and accepted 0 are $aq/(1+a)$ and
$(1-q)/(1+a)$. Their sum is at least $\min(q,1-q)\ge1/4$ and their
normalized ratio is $aq/(aq+1-q)=\sig(\ell+u)$.
\end{proof}
This is a two-coin/Bernoulli-race identity, established machinery
\cite{barker,bernoullirace}; the new use here is its physical interface and
composition with rare-event generation.

For a requested $0<\theta\le1/4$, set
\[
 J=4\lceil\log_2(8/\theta)\rceil,\qquad
 \beta=\theta/(16J).
\]
Cap the experiment after $J$ attempts, returning a fixed fallback bit on
exhaustion. The ideal tail is at most $(3/4)^J\le e^{-J/4}<\theta/8$.
Before each read hold for
$t_{\rm hold}\ge r_{\rm cal}^{-1}\log(1/\beta)$.
This makes the physical coin within conditional TV $\beta$ of $q$.
A finite dyadic branch selector has error at most $\beta$. The two
substitutions over at most $J$ attempts add at most $2J\beta=\theta/8$.

Let the numerical values of each of the actual fields $u,z$ be within
$\xi\le\theta/4$. The computed log shift then differs by at most $2\xi$;
since $\sig$ is $1/4$-Lipschitz, the resulting infinite-factory output
probability differs by at most $\xi/2\le\theta/8$. Acceptance remains at least
$1/4$ for every positive branch weight even with the perturbed shift.
Finite exponential/selector evaluation is included in the selector error.
Thus the entire capped output has conditional TV at most $\theta$, with
slack for the declared response/readout error allocations. Standard interval
exponentials and dyadic thresholds require bit cost polynomial in
$B_0+U+\log(J/\theta)$; their cost is not proportional to $e^{B_0+U}$.

The actual service bill per requested Gibbs update is bounded by
\begin{equation}\label{eq:gibbscost}
\begin{split}
 C_G(\theta)\le{}&
 C_{\rm context}+C_{\rm field}(\theta/4)+C_{\rm slew}+C_{\rm settle}\\
 &+J\!\left[r_{\rm cal}^{-1}\log(16J/\theta)+C_{\rm read}(\beta)
            +C_{\rm timer}(\Delta_G)\right]+C_{\rm arith}.
\end{split}
\end{equation}
The hold-timing resolution $\Delta_G$ is chosen so the actual hold is at
least the specified time; rounding upwards suffices if its excess time is
charged. This method requires no focal-state reset because its bias bound
holds uniformly conditionally on every earlier experimental history.

\begin{lemma}[Balanced log-odds measurement]\label{lem:odds}
At a balanced context, a finite number of physical readings,
$O(h^{-2}\log(1/\beta_0))$, and additive-field
calibration to $O(h)$ estimate $\ell_i(s)$ to absolute error $h$, with
conditional failure probability at most $\beta_0$. Every reading holds for
$O(r_{\rm cal}^{-1}\log(1/h))$. All costs are uniform over the previous
experimental transcript.
\end{lemma}
\begin{proof}
For $h\le1$, set $a=h/128$. Hold until each conditional mean differs from
the true $q_z$ by at most $a/2$. Martingale Hoeffding on
$M\ge2a^{-2}\log(2/\beta_0)$ observations bounds the empirical error by
$a$, including the conditional-mean bias, with probability at least
$1-\beta_0$. Clip the empirical probability to $[1/8,7/8]$. Since true
$q_z\in[1/4,3/4]$, clipping cannot increase its error. The logit derivative
on the clipping interval is at most $64/7<10$, so its error is at most
$10a<h/12$. Subtract the measured field $z$, whose actual value is known
to $h/4$, and evaluate the logarithm to $h/4$. The total is less than $h$.
For $h>1$ one may use the $h=1$ procedure. The hold and every calibration,
response, read and timer service are charged.
\end{proof}

No fine log-odds estimate is used in the Gibbs update of
Lemma~\ref{lem:gibbs}. Lemma~\ref{lem:odds} is used only when an actual
relative static probability needs a local weight-ratio evaluation.

\section{True operating updates from censored exits}

The actual operating clocks can change with the menu. We never infer them
from the clocks at balancing fields. Fix a requested coordinate, its
neighborhood, starting bit, and full operating environment $u$. Pin that
bit, release, and let its first-exit time be $T\sim\operatorname{Exp}(\lambda)$,
where $\lambda$ is the applicable rate in \eqref{eq:Q}. It obeys
$0<\lambda\le r^*$.

\begin{lemma}[Operating race and finite walk]\label{lem:walk}
Take $\rho=2r^*$, generate an independent
$E\sim\operatorname{Exp}(\rho)$, and censor the focal first-exit observation
at $E$. The bit $\one_{\{T<E\}}$ has probability
$p=\lambda/(\lambda+\rho)\le1/3$.
A random walk starting at 0, moving up 1 on a $p$-coin success and down 1
otherwise, ever hits $+1$ with probability
$p/(1-p)=\lambda/\rho$.
For $0<\theta\le1/4$, put
\[
 m=\lceil\log_2(8/\theta)\rceil,\qquad n=72m .
\]
Stop the walk at $+1$, at $-m$, or after $n$ steps; output 1 only if $+1$
has been reached. Its output probability differs from $\lambda/\rho$ by
less than $\theta/4$. If each implemented input race bit has conditional
TV error at most $\beta=\theta/(8n)$, the full output error is less than
$\theta/2$.
\end{lemma}
\begin{proof}
Independence of the two clocks gives
$\Prb(T<E)=\int_0^\infty\lambda e^{-(\lambda+\rho)t}\,dt$.
The infinite-walk hitting probability is the classical biased-walk solution
$r=p/(1-p)$, obtained by solving the bounded harmonic hitting equation or
optional stopping on finite intervals and taking a limit. Since $r\le1/2$,
the paths lost by absorption at $-m$ have probability at most the eventual
success probability starting there, namely $r^{m+1}\le2^{-m-1}<\theta/16$.

For the unlimited walk, $\E S_n\le-n/3$. A path not yet absorbed at either
boundary has $S_n>-m$. Since $n\ge6m$, Hoeffding for independent increments
in $[-1,1]$ gives
\begin{align*}
 \Prb(\hbox{not absorbed at }n)
 &\le\Prb(S_n-\E S_n>n/3-m)\\
 &\le\exp[-(n/3-m)^2/(2n)]\\
 &\le e^{-n/72}=e^{-m}\le2^{-m}\le\theta/8 .
\end{align*}
The reference walk uses independent ideal coins. Conditional kernel errors
of the physical implementation are coupled sequentially: the probability
any of at most $n$ races differs is at most $n\beta=\theta/8$.
Adding the three errors proves the result without assuming the physical
coins independent.
\end{proof}

The construction is a particular Bernoulli factory; the general analytic
Bernoulli-factory theory is established \cite{nacuperes}. The explicit
bounded walk, censored physical race and total error/cost accounting are
provided because they are needed by the claimed physical operation.

\subsection{A finite deadline and comparison device}

For each ideal race input at accuracy $\beta$, a sufficient finite recipe is
\begin{equation}\label{eq:timing}
 D=\rho^{-1}\lceil\log_2(32/\beta)\rceil,\qquad
 0<\Delta\le\beta/(256\rho).
\end{equation}
The ideal deadline tail is $e^{-\rho D}\le\beta/32$.
Use an independent uniform $b$-bit grid variable to implement the inverse
exponential map. Couple it to a true $U\in(0,1)$ by a uniform grid interval.
On $E=-\log U/\rho\le D$ the map's derivative is at most
$e^{\rho D}/\rho$. A sufficient uniform grid mesh is
\[
 2^{-b}\le\min\{\beta/256,\ \rho\Delta e^{-\rho D}/32\}.
\]
An endpoint cell intersecting $e^{-\rho D}$ has probability at most
$2^{-b}$ and is charged separately. The endpoint $U=0$ is never evaluated:
all grid cells below the cap threshold return the declared fallback.
On the other cells, finite interval evaluation and rounding produce a
scheduled deadline within $\Delta$ of the coupled ideal $E$. Choosing
interval width at most $\Delta/4$ and rounding to a clock grid of spacing
at most $\Delta/4$ supplies this bound. The needed bit length is
$O(\log(1/\beta)+\log(1/(\rho\Delta)))$; the displayed inequalities can be
checked with rational conservative exponential bounds.

Include actual deadline service error and timestamp/response errors in a
combined comparison error of at most $3\Delta$. The physical comparison can
then differ from the ideal event only if $|E-T|\le3\Delta$, or on a cap or
numerical failure. Conditional on $T$, the density of the independent $E$
is at most $\rho$, so the tie band has probability at most $6\rho\Delta$.
Use a fixed rule in ties. Thus cap, threshold cell, rounding and tie-band
errors sum to less than $\beta/8$ with the choices above; the remaining
budget up to $\beta$ covers declared read, missed-exit and service errors.
If these errors are deterministic accuracy bounds rather than random
failures, their contributions are just the tie-band bound. No total
variation approximation of a continuous real clock by a discrete variable
is asserted; only the event-comparison bit is approximated.

Each race ends by $D+O(\Delta)$, including a no-exit observation. It never
waits for an uncontrolled rare full dwell $1/\lambda$. But it uses a new
requested-state pin and release before every trial, even when that state
is rarely occupied naturally. The pin is a substantive resource.

\subsection{Full operating kernel and cost}

Choose a coordinate uniformly from $N$, request its actual operating
environment and neighborhood, and use Lemma~\ref{lem:walk} to flip it with
probability $\lambda_i/\rho$. The ideal kernel is exactly
$P=I+Q^u/(N\rho)$. A finite coordinate selector can use an error budget
$\theta/4$; the preceding factory has slack for the remaining budget.

For the actual source and target states,
$\pi^u(x)Q^u(x,x^{(i)})=\pi^u(x^{(i)})Q^u(x^{(i)},x)$, because the
same $r_i^u(s)$ multiplies the complementary logistic probabilities.
Therefore $P^*=P$. This route acquires both update interfaces only for the
declared reversible form. It does not construct the reversed kernel of an
arbitrary unknown nonreversible process.

A worst-case service cost for an operating update at error $\theta$ is
\begin{equation}\label{eq:flipcost}
\begin{split}
 C_P(\theta)\le{}&C_{\rm context}+C_{\rm field}(\xi_{\rm op})
                  +C_{\rm slew}+C_{\rm arith}\\
 &+n\big[C_{\rm pin}+C_{\rm settle}
       +C_{\rm event}(\beta)+C_{\rm timer}(\Delta)
       +D+O(\Delta)\big],
\end{split}
\end{equation}
with $m,n,\beta,D,\Delta$ as above and with any operating-environment
reproduction precision $\xi_{\rm op}$ required to produce the declared
actual $Q^u$. When faithful copies reproduce that environment exactly by
assumption, $C_{\rm field}(\xi_{\rm op})$ denotes that paid reproducible
setting. Otherwise the induced kernel difference is separately bounded
and allocated to $\theta$; an unknown arbitrary dependence of $r_i^u$
on the field does not permit a continuity inference from field accuracy.
Deployment must reproduce the same environment or carry an additional
proved law-error bound.

The number of pins is $O(\log(1/\theta))$, and the total censored dwell is
$O(\rho^{-1}\log^2(1/\theta))$ with the explicit conservative constants
above. Small rates need not be estimated: a small success probability is
implemented by a correspondingly unlikely success. This does not improve
the cost of \emph{estimating} a small rate relatively. It does not remove
the numerical $r_*^{-1}$ dependence of $g^{-1}$ or the cost of physically
setting a locally unlikely starting bit.

\section{Experimental rare static access without an unknown weight oracle}

\begin{lemma}[Finite experimental pinned sampler]\label{lem:pinned}
For any coordinate pinning $D$ and $0<\omega\le1/4$, a conditional-on-history
TV-$\omega$ draw from $\pi^u(\cdot\mid D)$ uses at most $NS$ Gibbs updates,
where
\[
 S=\left\lceil\alpha^{-1}\left\lceil\log_2(4N/\omega)\right\rceil\right\rceil,
 \qquad
 C_{\rm pinlaw}(\omega)\le
 NS\,C_G\bigl(\omega/(4NS)\bigr)+C_{\rm virtual}.
\]
The pin constraints here are virtual classical constraints. They are not a
claim that the whole physical system can be prepared in the rare pin event.
\end{lemma}
\begin{proof}
Start a virtual configuration in any pin-compatible state. Systematically
update its free coordinates, requesting a local cell at the current virtual
neighborhood each time. Pinning preserves the influence bounds.
For ideal Gibbs updates, if all coordinate disagreement probabilities before
a sweep are at most $d_0$, none exceeds $d_0$ during it and each updated
coordinate has disagreement at most $(1-\alpha)d_0$. A complete sweep gives
that bound everywhere. Coupling to a stationary copy gives TV at most
$N(1-\alpha)^S\le Ne^{-\alpha S}\le\omega/4$.
Substituting conditional-TV-$\omega/(4NS)$ physical updates costs at most
$\omega/4$ by sequential coupling. The result has slack for finite control
and any charged services. If all coordinates are pinned, return that state.
\end{proof}

The accepted heavy-prefix estimator used the numerical terminal weight
$W(x_D)$. It cannot be copied unchanged when $W$ is unknown. The following
identity closes precisely that gap.

\begin{lemma}[Experimental relative pin mass]\label{lem:pinmass}
For any nonempty coordinate-pinning event $D$, rational
$0<v,\beta_0\le1/4$, the physical access constructs a positive rational
$\widehat p_D$ with
\[
 \Prb\!\left(e^{-v}\pi^u(D)\le\widehat p_D\le e^v\pi^u(D)
       \,\middle|\,\hbox{prior transcript}\right)\ge1-\beta_0 .
\]
Its worst-case number of pinned draws is
$O(N^3v^{-2}\log(N/\beta_0))$; each draw has TV tolerance $O(v/N)$.
It additionally uses $O(N^3v^{-2}\log(N/\beta_0))$ balanced readings to
estimate local log odds, at field precision $O(v/N)$ and hold time
$O(r_{\rm cal}^{-1}\log(N/v))$. There is no inverse $\pi^u(D)$ cost.
\end{lemma}
\begin{proof}
Heavy-complete $D$ by visiting its free coordinates. At each stage estimate
the current conditional marginal and choose the empirically heavier bit,
then add that bit to the prefix. Independently in the algorithmic sense
(fresh future experiments, not a physical reset assumption), heavy-complete
the empty pin event. Denote their terminal states by $x_D,x_0$, and the
products of their true selected conditional probabilities by $H_D,H_0$.
The exact identities
\[
 H_D=\frac{\pi^u(x_D)}{\pi^u(D)},\qquad H_0=\pi^u(x_0)
\]
give
\begin{equation}\label{eq:unknownweight}
 \pi^u(D)=
   \frac{W^u(x_D)}{W^u(x_0)}\frac{H_0}{H_D}.
\end{equation}
When $D$ pins every coordinate, $H_D=1$ and the same identity applies.

Put $a=v/(32N)$. Estimate each of at most $2N$ marginal stages using
\[
 M\ge 2a^{-2}\log(8N/\beta_0)
\]
pinned draws, each at conditional TV at most $a/2$. Conditional mean bias
is at most $a/2$, and bounded martingale Hoeffding bounds the empirical
deviation around its conditional means by $a/2$ with failure probability
at most $\beta_0/(4N)$. A union bound gives a simultaneous event of
probability at least $1-\beta_0/2$.
On it, every empirically selected true heavy probability obeys
$h_j\ge1/2-a\ge3/8$ and its estimated value is at least $1/2$.
The log error of one factor is at most $(8/3)a$ by the mean-value theorem.
The two products thus contribute at most $(16/3)Na=v/6$ to the log error
in \eqref{eq:unknownweight}.

Connect $x_0$ to $x_D$ by changing their differing coordinates one at a
time, in any fixed order; there are at most $N$ steps. If step $j$ changes
coordinate $i_j$ of $x^{j-1}$, the log weight ratio is exactly
\begin{equation}\label{eq:pathodds}
 \log\frac{W^u(x_D)}{W^u(x_0)}
 =\sum_j(1-2x^{j-1}_{i_j})
        \bigl[\ell_{i_j}(x^{j-1}_{\partial i_j})+u_{i_j}\bigr].
\end{equation}
These are full local conditional log odds, not marginals under the rare
pin event. Estimate each bracket to accuracy $v/(4N)$ with
Lemma~\ref{lem:odds}, including the actual $u$ calibration and numerical
evaluation in that tolerance. Allocate failure $\beta_0/(2N)$ per term.
Their total log error is at most $v/4$ on a joint event of failure at most
$\beta_0/2$. The chosen contexts are adaptive, but the conditional
guarantees are uniform, so the union bound remains valid.

Compute a finite positive rational exponential of the resulting sum to
log-relative error at most $v/12$, and multiply by the rational empirical
heavy products. The total log error is at most
$v/6+v/4+v/12=v/2<v$, leaving slack. All selected empirical factors are at
least $1/2$. The absolute magnitude of \eqref{eq:pathodds} is at most
$N(B_0+U)$, so the rational output and computations have bit length
polynomial in $N(B_0+U)$, $\log(1/v)$, and the sample-count logarithms.
There are at most $2N$ heavy stages, each using $M=O(N^2v^{-2}
\log(N/\beta_0))$ pinned draws. At most $N$ balanced logit estimates have
the same order of individual readings. This proves the claimed counts.
\end{proof}

One explicit cost form, with the constants of the preceding proof, is
\begin{equation}\label{eq:masscost}
 \begin{split}
 C_{\rm mass}(v,\beta_0)\ \le{}&
 2NM\,C_{\rm pinlaw}(v/(64N))\\
 &+N\,C_{\rm odds}(v/(4N),\beta_0/(2N))
   +C_{\rm ratio-arith},
 \quad M=\left\lceil2(32N/v)^2\log(8N/\beta_0)\right\rceil .
 \end{split}
\end{equation}
Here $C_{\rm odds}$ includes every actual calibration, setting, slew,
settling, hold, read and timer call in Lemma~\ref{lem:odds}.

\subsection{From pin masses to disconnection Count and SampleB}

Let the physical failure event $B$ be all-terminal disconnection of the
declared graph, with bit 1 denoting an edge failure. The full failure rule
is supplied. The underlying graph is connected and has at least two
vertices, so $B$ is nonempty and the cut weights have a positive minimum.
The coarse rational bounds \eqref{eq:contrast}, with fixed
$c$, yield rational cut weights as in the accepted correlated construction.
For $C=\lceil2c\rceil$, that construction has two branches.
If its minimum cut weight $\lambda_{\rm cut}$ is below
$4\lceil\log_2(2v_{\rm graph})\rceil$, it supplies the explicit lower bound
$p_0=2^{-8C\lceil\log_2(2v_{\rm graph})\rceil}$ on $\pi^u(B)$; ordinary
Gibbs rejection then costs at most $v_{\rm graph}^{O(C)}$.
In the other branch, at target error $h$, it constructs at most
\[
 J_B\le64v_{\rm graph}^{4C+6}/h
\]
cut events and omits at most $(h/16)\pi^u(B)$ in mass on its declared
enumeration event.

Use Lemma~\ref{lem:pinmass} to estimate each cut-pinning probability $p_j$
by a positive rational $\phi_j$ with
$e^{-s}p_j\le\phi_j\le e^sp_j$. Draw index $j$ proportional to $\phi_j$,
then draw from the pinned law with Lemma~\ref{lem:pinned}. For exact
component draws, the proposal density is
\[
 \nu(x)=\frac{\pi^u(x)}{\sum_j\phi_j}
       \sum_j\frac{\phi_j}{p_j}\one_{E_j}(x).
\]
With $d(x)=\sum_j\one_{E_j}(x)$, the ratio of the sum in this expression
to $d(x)$ lies in $[e^{-s},e^s]$.
The reciprocal statistic $1/d(x)\in[1/J_B,1]$ therefore gives relative
union counting by a bounded-mean estimator; accepting with probability
$1/d(x)$ gives the union-conditioned law with a multiplicative density
error $e^{\pm2s}$. This is the exact place where relative pin mass,
rather than knowledge of $W$, is needed.

For counting at relative error $h$, the accepted schedule uses
$s=h/16$, joint pin-mass confidence error at most one quarter of the
requested count failure probability, and at most
\[
 \left\lceil2048J_Bh^{-2}\left\lceil\log_2(8/\beta_0)\right\rceil\right\rceil
\]
union proposals.
For a conditional draw at TV error $h$, it uses a fresh list, total
enumeration/mass failure at most $h/8$, $s=h/32$, and at most
$\lceil2J_B\lceil\log_2(16/h)\rceil\rceil$ proposals. It always returns
a state in $B$, using the known all-failed state as the finite-cap fallback.
Finite categorical and reciprocal arithmetic are the accepted finite
rational procedures. All requested relative pin accuracies and confidences
are thus explicit functions of $h,\beta_0,J_B$.

Dependent physical calls are handled by conditional means, rather than
requiring every sample to succeed on a shared event.
Condition on the good retained-list and mass-estimate event. For the
count statistic $Y=1/d(X)$, every subsequent conditional mean differs
from the ideal mixture mean by at most the proposal-TV and rounding
budget $h/(64J_B)$. The ideal mean is at least $1/J_B$.
Conditional bounded-variable exponential-moment concentration, proved in
Lemma~\ref{lem:means}, therefore gives the same relative concentration
schedule, with confidence entering only logarithmically.
For SampleB, every proposal conditional on every previous history is
within its allocated TV of the same ideal proposal. Its accepted mark law
is uniformly close to the same union-conditioned law by
Lemma~\ref{lem:accept}. A mixture over the possible first accepted
histories retains that bound, and the cap has the same exponential tail.
Thus the accepted static tolerance schedule is valid with fresh
conditional-on-history physical calls, without an iid assumption.
Running it at half the requested error supplies slack for conservative
constants. Pin-mass confidence failures for SampleB are included in its
unconditional TV budget; no union over every future SampleB invocation
is demanded.

Consequently the exact interfaces needed by the accepted first-event
sampler are now physical finite operations:
\[
 \operatorname{Count}(h,\beta_0),\qquad
 \operatorname{SampleB}(h),
\]
with conditional-on-history confidence/TV guarantees for the true law.
A fresh SampleB call rebuilds its random estimates. Reusing a shared
random table instead requires an explicitly allocated outer confidence
event; it is not an independent sample by declaration.

\subsection{Optional fitted-static alternative}

A distinct implementation may use the accepted active learner to construct
a rational static weight table $\widehat W$ satisfying, on one shared
confidence event,
$e^{-E}\le\widehat W(x)/W(x)\le e^E$ for every $x$.
Then
\[
 e^{-2E}\le\frac{\widehat\pi(B)}{\pi(B)}\le e^{2E},
 \qquad
 \TV\bigl(\widehat\pi(\cdot\mid B),\pi(\cdot\mid B)\bigr)\le\tanh E .
\]
Indeed normalization on the whole space or on $B$ changes a pointwise
weight ratio by at most another factor $e^E$, and a likelihood ratio in
$[e^{-2E},e^{2E}]$ has TV at most $\tanh E$.
Choose $E$ to spend at most half the requested count/sample error;
the known-model static construction supplies the other half. Its pinned
mixing and fixed contrast must be certified for the fitted weights too.
True physical $P=P^*$ updates still use the operating-exit factory, so no
fine kinetic table is needed.
This hybrid can amortize an acquired static table over many records.
The on-demand weight-ratio route is an alternative, not a universal
cost improvement over that hybrid.


\section{Conditional errors, finite confidence and the complete operation}
The local cells can retain arbitrary previous states. What the proofs
supply is a response kernel uniformly close to the intended one,
conditional on every preceding experimental transcript. They do not
assert that reused physical cells are independent.

\begin{lemma}[Finite adaptive substitution]\label{lem:rootAdaptive}
If a finite controller makes at most $D$ adaptive calls, and each
conditional response differs from its ideal kernel by at most
$\theta_j$, with $\sum_{j\le D}\theta_j\le a$ deterministically, its
whole transcript differs from the ideal one by at most $a$ in TV.
\end{lemma}
\begin{proof}
Couple the classical randomness. Before the first mismatch the requested
inputs agree. Couple the next response maximally, and sum the conditional
first-mismatch probabilities. The claim holds for predictable allocations
and nested calls with finite caps.
\end{proof}

A naive choice to implement every \emph{whole static sampler} at accuracy
$\delta/D$ would unnecessarily worsen its accuracy cost. There are two
valid alternatives: flatten to primitive physical coins whose accuracy
cost is logarithmic, paying their actual metrology services, or use
conditional-mean concentration directly. The latter gives a transparent
confidence interface.

\begin{lemma}[Bounded conditional-mean calibration]\label{lem:rootChernoff}\label{lem:means}
Let $Z_1,\ldots,Z_n\in[0,1]$ obey, for every prior transcript,
$\E[Z_j\mid\mathcal F_{j-1}]\in[\mu-b_0,\mu+b_0]$.
Suppose $\mu\ge q_0>0$, $0<\xi\le1/2$ and
$b_0\le\xi q_0/16$. Then
\[
n\ge\frac{32}{\xi^2q_0}\log\frac2\delta
\quad\Longrightarrow\quad
\Prb(|n^{-1}\sum_j Z_j-\mu|>\xi\mu)\le\delta.
\]
\end{lemma}
\begin{proof}
Convexity on $[0,1]$ gives
$\E[e^{tZ_j}\mid\mathcal F_{j-1}]
\le\exp((\mu+b_0)(e^t-1))$ for $t>0$.
For $t<0$ the same bound uses $\mu-b_0$.
Iterating conditional expectations gives the usual Bernoulli upper and
lower exponential bounds with these two fixed means; independence is
not needed. The distance between either target threshold
$\mu(1\pm\xi)$ and the appropriate comparison mean is at least
$15\xi\mu/16$. The elementary multiplicative Chernoff inequalities
then bound the sum of the two tails by
$2e^{-nq_0\xi^2/12}$; the displayed constant is conservative.
\end{proof}

\begin{lemma}[Conditional accepted records]\label{lem:accept}
An ideal attempt outputs rejection or an accepted mark, with success
probability $q\ge q_0>0$ and conditional mark law $\nu$. Suppose that
after every preceding history the actual combined rejection/mark law
is within TV $a\le q_0/2$ of this ideal law. Its conditional successful
mark is within $2a/q_0$ of $\nu$. Capping at $J$ attempts and returning
an arbitrary fallback mark on exhaustion has output TV at most
$2a/q_0+e^{-q_0J/2}$ from $\nu$.
\end{lemma}
\begin{proof}
The success mass changes by at most $a$. Normalizing the two accepted
submeasures and adding and subtracting their common-denominator form
gives the stated $2a/q_0$ bound. This bound is uniform over every
pre-success history and hence over their mixture. Each conditional
success probability is at least $q_0/2$, so exhaustion has probability
at most $(1-q_0/2)^J\le e^{-q_0J/2}$. No independence is assumed.
\end{proof}

For Palm calibration, $Z$ is the rounded reciprocal visit count. A
proposal's conditional TV error bounds its conditional mean bias, and
the accepted floor is $q_0=g/8$. Choose this bias at most
$\xi q_0/16$; the number of proposals remains
$O(g^{-1}\xi^{-2}\log(1/\delta))$. Count's reciprocal-overlap
statistics similarly have a polynomial floor $1/J_B$. Heavy-prefix
marginal estimation already uses martingale Hoeffding. These facts
preserve logarithmic confidence dependence in the abstract query counts.
An apparatus service such as $C_{\rm field}(\theta)$ may have a much
worse precision dependence; its physical cost is still paid.

\begin{theorem}[Acquired joint-event generation from local queries]
Under the faithful-cell model, calibrated and reproducible environments,
finite physical services, certified pinned margin and fixed contrast
above, the all-terminal disconnection first-event record admits the
following finite operation. With prescribed shared confidence
$1-\delta$, a compiled generator's next record, conditional on every
previous experimental transcript, is within prescribed TV error
$\varepsilon$ of the true quantized record from a supplied initial law.
The initial-law approximation is included in the error budget. The
generator uses fresh local experiments on every draw, without a fine
table of operating rates. A fitted static table is optional.

Compilation and per-record query counts are finite and polynomial in
the displayed static interfaces at fixed contrast exponent, numerical
inverse mixing margins, $\log(1/\zeta)$, the value of the requested
window, clock digits, inverse error and log inverse confidence.
Degree dependence is that of the chosen coarse certificate and the
priced neighborhood interface. There is no inverse global-rarity or
mean-lifetime trajectory factor. Every apparatus service is evaluated
at its actual requested accuracy and its cost is added.
\end{theorem}
\begin{proof}
The three preceding constructions supply Count, SampleB and $P=P^*$
for one and the same true invariant law and operating generator.
Their conditional errors can be budgeted through finite adaptive
substitution for records and through the conditional concentration
lemma for calibration. The accepted first-event compiler then applies,
including its rare/direct branch, arbitrary-start correction, finite
clock, numerical guards and capped loops. Its mathematical Palm and
mixing arguments do not require the unknown physical rates to be
rational; the physical kernel calls replace their finite numerical
evaluation. Allocate, for example, half the final record error to its
finite construction and the remainder to physical substitutions and
initialization, with the corresponding independent confidence budgets
for structural and static/shared calibration events. Each nested
budget has a finite call cap and can be refined before execution.
All quantities are common to the true physical law; no proxy stationary
law is inserted without its own error bound.
\end{proof}

For explicit accounting, let $M_{a,G},M_{a,P},M_{a,O}$ be the finite
numbers of Gibbs updates, operating updates and balanced log-odds
estimates after expanding the accepted static/sampler schedules, with
$a=0$ for compilation and $a=1$ for a record. Then
\begin{equation}
\begin{split}
C_a\le C_{a,\rm classical}
&+\sum_{j=1}^{M_{a,G}}C_G(\theta_{a,j})
+\sum_{j=1}^{M_{a,P}}C_P(\theta_{a,j})\\
&+\sum_{j=1}^{M_{a,O}}C_{\rm odds}(h_{a,j},\beta_{a,j}).
\end{split}
\label{eq:physicalLedger}
\end{equation}
Coarse acquisition, shared metrology, available hardware, propagation,
stored configurations/cut lists/counters and chosen-intervention
deployment/preparation are additional explicit entries. Parallelism
changes wall time only by paying the corresponding cells and services.
For $K$ menu items and $n$ records per loss estimate, the decision pays
the $K$ compilations and $Kn$ per-record costs. Conditional Hoeffding
and the same empirical-minimizer argument as \eqref{eq:regret} give
the bounded-loss decision guarantee without an independence assumption
on actual reused cells.

For physical equilibrium-safe missions, keep the certified safe-mass
and finite preparation conditions of Section~\ref{sec:acquired}; include
preparation error and abort loss. With true experimental static draws,
there is no fitted-weight $\tanh(E_u)$ term. The actual chosen operating
environment must be reproduced. Approximate field reproducibility alone
does not control an arbitrarily discontinuous unknown clock response.
Either faithful reproduction is supplied, or a separate rate/law bound
must be acquired and charged. This is one of the physical premises
that the experiments do not eliminate.


\section{What the representation changes}
\subsection{An explicit failure of full-lifetime information accounting}
Let $X$ be a two-state component, initially safe, with safe-to-failed
rate $\lambda$ and repair rate one. Independently let $Y$ be a symmetric
two-state chain with switch rate $c$, initially uniform. The event is the
first failure of $X$. Compare it with a model having the same $X$ and
switch rate $\widehat c$ for $Y$. Both event CDFs are identical and both
terminal $Y$ laws are uniform.

Write $j=c\log(c/\widehat c)-c+\widehat c$. The KL divergence of the
entire $Y$ trajectory up to the first failure is $j/\lambda$.
For a retained physical window $h$, however, condition on $\tau$ and
use stationarity of $Y$: the retained segment has length at most $h$,
common uniform initial law and KL at most $hj$. Thus the joint time,
terminal state and window obey TV at most $\sqrt{hj/2}$, independently
of $\lambda$. Long stretches of unreported ordinary behavior can carry
arbitrarily much model-discriminating information while contributing no
corresponding uncertainty to the retained first-event question.

This observation alone is not an algorithm for interacting systems.
The proposal construction in Section~\ref{sec:bridge} supplies the
computable interacting-system replacement: local divergences are paid
on a short forward/reverse proposal, its successful record is controlled
through accepted submeasures, and only a relative block-frequency error
is paid by the binary waiting clock.

\subsection{Comparison with strong applicable methods}
Goal-oriented path-space uncertainty is established machinery.
Birrell, Katsoulakis and Rey-Bellet explicitly treat stopping-time
distributions and general stopped observables using appropriate
relative-entropy objects \cite{birrell}. R\'enyi-divergence methods
give rare-event sensitivity bounds that can remain useful where ordinary
event-probability estimates fail \cite{renyi}. The contribution here is
the explicit local-evidence-to-marked-proposal bound, the nonrare
conditioning floor, and its finite acquisition/generation/decision
composition. We do not claim that all preceding uncertainty methods pay
inverse rarity, or that the use of an information distance is new.

Reverse-time sequential Monte Carlo already targets stopped rare paths
\cite{reverse}. Blanchet obtains strongly efficient overflow-path
sampling for Jackson networks, with work proportional to the path length
in the stated network family \cite{blanchet}. Parallel replica methods
use quasi-stationary structure to accelerate metastable trajectories
\cite{parrep}. These are substantial existing operations, and exact
finite-state solvers can be superior for small systems. The present
construction adds a finite joint-TV and experimental-interface certificate
for its declared succinct interacting model class; it is not a universal
replacement for those algorithms. Our certified gap and static-access
premises can fail in systems where a specialized rare-path method works.

Bernoulli factories and sampling from unknown probability coins are also
established \cite{huber,bernoullirace}. Our release race and bounded random-walk
factory are explicit elementary realizations for the particular physical
rate interface, not claims to invent the factory principle. Active
perturbations and nonequilibrium data can already make spin-model
inference much easier \cite{activefields,dutt}. The new operational
endpoint joins controlled local evidence to a finite rare first-event
record and its intervention risk, while pricing actual clocks, fields,
pins, preparation and fresh-query use.

A targeted search also found 2026 primary abstracts by Sun, Moody and
Barthel (arXiv:2608.16171) and Barthel, Sun and Peng
(arXiv:2608.17511) describing conditional-path Monte Carlo with
rejection-free spacetime cluster updates. Their full texts could not be
retrieved in this session. We record this relevant newer alternative
without asserting a theorem-level separation from it or using its
claims as proof premises.

\section{Executed evidence and its limits}
\subsection{An exact-rational uncertainty calculator}
The accompanying \texttt{certificate.py} calculates conservative
upper bounds with rational arithmetic. For $0\le a<1$, it replaces
$e^a-1$ by $a/(1-a)$; it bounds $e^{-gb}/\zeta$ by
$2^{-\lfloor gb\rfloor}/\zeta$ and rounds square roots upward using
integer square roots with verified rational endpoints. For
$0\le d<1$, it bounds $e^d-1$ by $d/(1-d)$. Exact rational
certificates are saved separately from decimal display values.
This is not a general Gibbs counting or first-event simulation engine.

The three-state example of Section 1 was executed with
$r=0.4$, $\widehat r=r(1+10^{-12})$ and
$\widehat\lambda=\lambda(1+10^{-12})$.
The weights $(1,\lambda r,\lambda(1-r))$ and their counterparts
satisfy $E=a=2\cdot10^{-12}$. The common choices
$R=4$, $g=1/4$, $\zeta=\lambda/8$, $p_+=2\lambda$ are valid.
For $H=1$ and schedule parameter $\varepsilon=0.2$, the calculation gave:

\begin{center}
\begin{tabular}{rrrr}
\toprule
Mean first time & Buffer $b$ & Proposal length $L$ & Certified model TV \\
\midrule
$10^{12}$ & 212 & 4,403,200 & $<0.004909$ \\
$10^{30}$ & 448 & 9,236,480 & $<0.006221$ \\
$10^{100}$ & 1,380 & 28,323,840 & $<0.005368$ \\
\bottomrule
\end{tabular}
\end{center}

In this example the pre-event state is constant, so the precursor record
is determined by the time and terminal mark. An independent analytic
density integral gives approximately $4.8405290\cdot10^{-13}$ for the
actual time-and-mark TV, the same for all three $\lambda$ values.
That value is a high-precision floating diagnostic, not the rational
certificate. The gap between it and the certificate shows how loose the
generic sufficient schedule is. The ordinary exact three-state solver
is far better for this example; this table establishes neither speedup
nor practical acquisition cost. It checks the new bound and its absence
of a numerical lifetime factor. The general ordered-precursor assertion
rests on the proof, not this trivial precursor execution.

An initial attempted $\lambda=10^{-6}$ instance failed the calculator's
rare-branch condition because its prescribed block was too long for
$Lp\le1$. The guard rejected it; no certificate was issued. The common
branch theorem instead directs such inputs to the finite direct branch.
The diagnostic subsequently used only the displayed rare instances.
This failed attempt is retained as a scope check.

\subsection{Independent constructions retained}
The separate decision-route execution checks its finite local factories
and static weight-ratio identity: 10 exact rational bounded-walk cases,
20 balanced-Gibbs factory cases, and 4 independently enumerated pin-mass
identities passed. The cases include operating rate $2^{-100}$ and a
pin event of probability approximately $1.45\cdot10^{-11}$. These are
exact arithmetic checks of finite formulas, not observations of those
rare outcomes in a physical system.

The service example is explicit. For a single operating update at
$\theta=2^{-20}$ and $\rho=2$, the prescribed cap is 1,656 pins, a
deadline of 19.5 model time units per race (20 if rounded to whole
units), and aggregate capped dwell 32,292 (33,120 when rounded).
The comparison timing allowance is at most
$1.406\cdot10^{-13}$ model time units, with a sufficient 109-bit
deadline-uniform draw. Context setup, pins, settling, detector service,
reference-clock metrology and computation are added. This substantial
bill is sufficient, not necessary, and demonstrates no economical
device. A full event forecast includes many such calls and the static
operations.

The complete record states which primitives were simulated and charges
the finite call caps. No physical
local cell, detector or controlled network was built. Neither that
execution nor the present calculator implements the full correlated
static cut algorithm on a large network.

An independent foundational route produced an explicit finite rational
compiler from two-counter programs to globally certified
$1/2$-contractions of $\mathbb R^4$ with known unique fixed point zero.
Exact finite arrival at zero and one fixed homogeneous sign event encode
halting. Its finite compiler and exact arithmetic checks were executed.
The late distinguishing signal shrinks exponentially; positive absolute
event margins have a finite logarithmic horizon. General piecewise-affine
reachability undecidability is prior machinery and scaling may already
imply this consequence. It is retained as an exact-versus-finite-resolution
boundary, without a priority or robust physical-computation claim. Its
full proof, primary comparisons and precision bill are in
\texttt{next5\_fixedlaw}.

The independent mechanical route asks whether active forces can replace
an absent relaxation certificate when inferring zero-frequency transport.
For its equilibrium linear family, true transport is $K=M$ versus
reference $K=1$, while every causal force controller with
$\int_0^T u_t^2dt\le E$ satisfies
\[
D_{\rm KL}(P_{a,M}^{\rm record}\Vert P_0^{\rm record})
\le a(1+T)+a^2ET^2,
\quad 0<a\le1/4,\quad M\ge2.
\]
Sending the visible coupling $a$ toward zero hides arbitrarily large
transport from any fixed finite time/effort budget. Its separate
finite-slew, specified-pair pulse reaches the cubic active-visibility
scale $U^2a^2T^3$. The proof and finite-meter/work-reserve diagnostic
retain their supplied equilibrium initialization, force port and exact
model-order qualification. The reference has minimal order one;
the alternatives have order two. No fixed-reference stopping
impossibility is asserted for an independently certified exact-order-two
class. These results do not alter the declared-model premises of the
acquired-event construction.

\section{Preservation and completion ledger}
The complete preceding first-event package contains 182 files, including
its preserved predecessors. The new package retains those originals
byte for byte, with the baseline and final SHA-256 comparison. Earlier
quantum acquisition/control/calibration, effusion, reliability,
transport, passive-memory, fixed-law, active acquisition/prevention and
fixed-mechanics results keep their original hypotheses and records.

The new completed objects are the marked-law comparison and its local
envelopes; the common finite branch and acquired uncertainty/decision
composition; the experiment-fed static and dynamic interfaces with
their finite caps; and the two independent boundary constructions.
The package also keeps exact source files, programs, execution outputs,
resource accounting, primary-source links, and unsuccessful approaches.
No original proof is replaced by a relaxed summary.

The substantive outstanding engineering task is economical realization
in a specified physical system: faithful cells, calibrated strong fields,
low-cost pins, fast detectors and strict mixing margins must actually
be available. The theorem does not claim that such hardware or an
economical general implementation has been demonstrated. It resolves
the declared scientific question: local controlled access can support a
joint rare-event forecast and a quantitative permitted intervention
decision without observing a global failure or simulating its full
lifetime, with all remaining access and numerical costs exposed.



\begin{thebibliography}{20}
\bibitem{governing}
\emph{First-event research package and governing predecessors}, fixed revision
\texttt{719e201332992002528c221c49f80f69431485f6}.
\url{https://github.com/harjassand/astra-research-knowledge/tree/719e201332992002528c221c49f80f69431485f6/research/returns/2026-11-05-p3-first-event}.
Original proofs, finite compiler, executions, qualifications and earlier
packages are reproduced without modification in the preservation tree.

\bibitem{birrell}
J. Birrell, M. A. Katsoulakis and L. Rey-Bellet,
\emph{Quantification of model uncertainty on path-space via goal-oriented
relative entropy}, ESAIM: Mathematical Modelling and Numerical Analysis
55 (2021), 131--169. DOI: 10.1051/m2an/2020070.
\url{https://arxiv.org/abs/1906.09282}.

\bibitem{renyi}
P. Dupuis, M. A. Katsoulakis, Y. Pantazis and L. Rey-Bellet,
\emph{Sensitivity analysis for rare events based on R\'enyi divergence},
Annals of Applied Probability 30 (2020), 1507--1533.
DOI: 10.1214/19-AAP1468. \url{https://arxiv.org/abs/1805.06917}.

\bibitem{huber}
M. Huber, \emph{Optimal linear Bernoulli factories for small mean problems}.
\url{https://arxiv.org/abs/1507.00843}.
See also \emph{Nearly optimal Bernoulli factories for linear functions},
DOI: 10.1017/S0963548315000371.

\bibitem{bernoullirace}
S. Dughmi, J. Hartline, R. Kleinberg and R. Niazadeh,
\emph{Bernoulli factories and black-box reductions in mechanism design},
STOC 2017. Author manuscript:
\url{https://www.cs.cornell.edu/~rdk/papers/stoc17-rad.pdf}.

\bibitem{barker}
F. B. Gon\c{c}alves, K. \L{}atuszy\'nski and G. O. Roberts,
\emph{Barker's algorithm for Bayesian inference with intractable likelihoods}
(2017). \url{https://arxiv.org/abs/1709.07710}.

\bibitem{nacuperes}
S. Nacu and Y. Peres, \emph{Fast simulation of new coins from old},
Annals of Applied Probability 15 (2005), 93--115.
DOI: 10.1214/105051604000000549.
\url{https://arxiv.org/abs/math/0309222}.

\bibitem{reverse}
J. Koskela, D. Span\`o and P. A. Jenkins,
\emph{Inference and rare event simulation for stopped Markov processes via
reverse-time sequential Monte Carlo}, Statistics and Computing 28 (2018).
DOI: 10.1007/s11222-017-9722-1.
\url{https://wrap.warwick.ac.uk/id/eprint/84735/3/WRAP_st-160317-wrap_-reverse_time_smc.pdf}.

\bibitem{blanchet}
J. Blanchet, \emph{Optimal sampling of overflow paths in Jackson
networks}, Mathematics of Operations Research
38 (2013), 698--719. DOI: 10.1287/moor.2013.0586.
Author version: \url{https://web.stanford.edu/~jblanche/papers/RareJackson7.pdf}.

\bibitem{parrep}
D. Aristoff, T. Leli\`evre and G. Simpson,
\emph{The parallel replica method for simulating long trajectories of
Markov chains}, Applied Mathematics Research eXpress (2014).
DOI: 10.1093/amrx/abu005.
\url{https://www.math.colostate.edu/~aristoff/parrep_dynamics.pdf}.

\bibitem{activefields}
J. Jiang, D. A. Sivak and M. Thomson,
\emph{Active learning of spin network models},
\url{https://arxiv.org/abs/1903.10474}.

\bibitem{dutt}
A. Dutt, A. Lokhov, M. D. Vuffray and S. Misra,
\emph{Exponential reduction in sample complexity with learning of Ising
model dynamics},
Proceedings of Machine Learning Research 139 (2021).
\url{https://proceedings.mlr.press/v139/dutt21a.html}.
\end{thebibliography}
\end{document}
